% Book 4, Year 2, chapter 4 -- Templates and generic programming.
% Provenance (never printed): sources/book4/BRIEFS.md, chapter 4.
\chapter{Templates and Generic Programming}\label{ch:b2:cpp-templates}

Last year's dynamic array held \code{int}s. The weather station needs one of \code{double}s, the library one of
titles. C offers two ways out, and the program below shows both: copy the code and change the type, or write it once
for blocks of bytes, with \code{void *} and an element size.

\omcode{code/bachelor-2/04-cpp-templates/listings/int_vec.c}{21}{38}{The array once more, for any element size, with
\code{void *}}\label{lst:b2:cpp-templates:int-vec}

\omoutput{code/bachelor-2/04-cpp-templates/outputs/int_vec.txt}{An array of \code{int}, written for \code{int}, and
an array of \code{double}, written for bytes.}\label{out:b2:cpp-templates:int-vec}

The first way multiplies the code: a bug fixed in the \code{int} version stays in the \code{double} one. The second
keeps one copy but throws the types away: \code{any_vec_push} accepts the address of anything, copies bytes one by
one with a loop the compiler cannot specialise, and every reader must cast \code{data} back to the right type, which
nothing checks. The C library's \code{qsort} is written the second way. C++ offers a third: write the code once, with
the type as a parameter, and let the compiler write each typed copy.

\begin{recall}
\cref{ch:b2:cpp-copy-move}: a class that owns memory follows the \omterm{def:b2:cpp-copy-move:rules}{rule of five}, moves are cheap and \code{noexcept},
and a growing array makes fewer than $2(n-1)$ transfers when it doubles. Book 3 sorted with \code{qsort}, which
receives the array as \code{void *}, the element size, and a comparison function called through a pointer.
\end{recall}

\section{Function templates}\label{sec:b2:cpp-templates:functions}

\begin{definition}[Function template, class template, template parameter]\label{def:b2:cpp-templates:template}
A \emph{template}\index{template (C++)} is a declaration preceded by \code{template <typename T, ...>}, whose
\emph{template parameters}\index{template parameter} stand for types (\code{typename T}) or for constant values (such
as \code{std::size_t N}). A \emph{function template}\index{function template} describes a family of functions, a
\emph{class template}\index{class template} a family of classes, one member of the family for each choice of the
parameters.
\end{definition}

\omcode{code/bachelor-2/04-cpp-templates/listings/max_template.cpp}{5}{20}{One \omterm{def:b2:cpp-templates:template}{function template}, three
types}\label{lst:b2:cpp-templates:max}

\omoutput{code/bachelor-2/04-cpp-templates/outputs/max_template.txt}{Each call used the version for its
type.}\label{out:b2:cpp-templates:max}

\begin{definition}[Template argument deduction]\label{def:b2:cpp-templates:deduction}
At a call of a \omterm{def:b2:cpp-templates:template}{function template} with no explicit arguments, the compiler finds the \omterm{def:b2:cpp-templates:template}{template} arguments by
\emph{template argument deduction}\index{template argument deduction}: it matches each parameter type against the
type of the corresponding argument. Each \omterm{def:b2:cpp-templates:template}{template parameter} must be deduced to a single type from all the arguments
in which it appears; a parameter that appears in no function parameter (only in the return type, say) cannot be
deduced, and must be given explicitly, as in \code{larger<double>(3, 7.5)}.
\end{definition}

Deduction does not convert: \code{larger(3, 7.5)} would make \code{T} both \code{int} and \code{double}, and the call
does not compile. Given explicitly, \code{T = double} turns the call into an ordinary one, where 3 converts to
\code{double}.

\section{Class templates}\label{sec:b2:cpp-templates:classes}

The growable array of the opening, written once for every type, is a \omterm{def:b2:cpp-templates:template}{class template}. Unlike the C version, it must
also respect what a C++ element needs: a \code{std::string} cannot be copied byte by byte, and its \omterm{def:b2:cpp-classes-raii:constructor}{constructor} and
\omterm{def:b2:cpp-classes-raii:constructor}{destructor} must run. The array therefore allocates \emph{raw} memory (\code{::operator new}, the C++ form of
\code{malloc}), builds each element in it with a \emph{placement new}, \code{new (address) T(...)}, and destroys it
with an explicit call to its \omterm{def:b2:cpp-classes-raii:constructor}{destructor}.

\omcodes{code/bachelor-2/04-cpp-templates/listings/vec.hpp}{8-36,72-75}{\code{Vec<T>}: the special members, \code{push_back} and
the data}\label{lst:b2:cpp-templates:vec}

\omcode{code/bachelor-2/04-cpp-templates/listings/vec.hpp}{37}{63}{\code{reserve} and
\code{clear}}\label{lst:b2:cpp-templates:vec-reserve}

\code{reserve} builds the elements in the new block before destroying the old ones. It moves them when their \omterm{def:b2:cpp-copy-move:move}{move
constructor} cannot throw and copies them otherwise (\code{std::move_if_noexcept}), so that if a copy throws the old
block is still complete: the strong guarantee of \cref{ch:b2:cpp-copy-move}, which \code{tests/vec_test.cpp} checks
with an element whose fourth copy throws. The remaining members are one-liners: \code{operator[]}, \code{at}, which
checks its index, \code{size}, \code{capacity}, and \code{begin} and \code{end}, which make the range-\code{for} work.

\begin{omfigure}
\begin{tikzpicture}[font=\footnotesize]
\node[anchor=west, font=\footnotesize\bfseries] at (-0.3,1.75) {a \texttt{Vec<std::string>} object (24 bytes)};
\node[cell, minimum width=1.5cm] (d) at (0.5,0.5) {$\bullet$};
\node[cell, minimum width=1.5cm, right=0mm of d] (s) {2};
\node[cell, minimum width=1.5cm, right=0mm of s] (c) {4};
\foreach \n/\l in {d/data\_, s/size\_, c/cap\_} \node[font=\scriptsize\ttfamily, below=0.5mm of \n] {\l};
\node[cell, minimum width=2.4cm, fill=omDef!10, font=\scriptsize\ttfamily] (e0) at (6.4,0.5) {Mill Lane};
\node[cell, minimum width=2.4cm, fill=omDef!10, font=\scriptsize\ttfamily, right=0mm of e0] (e1) {Quay Street};
\node[cell, minimum width=2.4cm, fill=black!5, font=\scriptsize, right=0mm of e1] (e2) {raw};
\node[cell, minimum width=2.4cm, fill=black!5, font=\scriptsize, right=0mm of e2] (e3) {raw};
\draw[ptr] (d.center) .. controls (1.2,1.35) and (4.4,1.35) .. (e0.north west);
\draw[decorate, decoration={brace, mirror, amplitude=4pt}] (e0.south west) -- (e1.south east)
  node[midway, below=5pt, font=\scriptsize] {\texttt{size\_} elements, constructed};
\draw[decorate, decoration={brace, mirror, amplitude=4pt}] (e2.south west) -- (e3.south east)
  node[midway, below=5pt, font=\scriptsize] {room for \texttt{cap\_ - size\_} more, no object};
\end{tikzpicture}
\par\smallskip{\small The layout of every \code{Vec<T>}: three members of 24 bytes in all, whatever \code{T}, and a
heap block of \code{cap_} slots of \code{sizeof(T)} bytes, the first \code{size_} of which hold constructed objects.
The \omterm{def:b2:cpp-classes-raii:constructor}{destructor} destroys exactly those, then frees the block.}
\end{omfigure}

\omcode{code/bachelor-2/04-cpp-templates/listings/vec_demo.cpp}{6}{25}{One \omterm{def:b2:cpp-templates:template}{class template}, three
classes}\label{lst:b2:cpp-templates:vec-demo}

\omoutput{code/bachelor-2/04-cpp-templates/outputs/vec_demo.txt}{Integers, doubles and strings in the same
code.}\label{out:b2:cpp-templates:vec-demo}

\begin{definition}[Generic programming]\label{def:b2:cpp-templates:generic}
\emph{Generic programming}\index{generic programming} writes algorithms and data structures once, in terms of the
operations they need from their element types (copy, compare, add), so that they apply to every type providing
those operations, with no loss of type checking or of speed against a version written for one type.
\end{definition}

\section{Instantiation}\label{sec:b2:cpp-templates:instantiation}

A \omterm{def:b2:cpp-templates:template}{template} is not code. It becomes code when it is used.

\begin{definition}[Template instantiation]\label{def:b2:cpp-templates:instantiation}
\emph{Template instantiation}\index{template instantiation} is the generation, by the compiler, of a function or a
class from a \omterm{def:b2:cpp-templates:template}{template} and a list of \omterm{def:b2:cpp-templates:template}{template} arguments, by substituting the arguments for the parameters. A use
such as \code{Vec<int>} or \code{larger(3, 7)} causes the instantiation; a \omterm{def:b2:cpp-classes-raii:class}{member function} of a \omterm{def:b2:cpp-templates:template}{class template} is
instantiated only if it is used.
\end{definition}

\omcode{code/bachelor-2/04-cpp-templates/listings/instantiate.sh}{2}{4}{Compiling the demo and listing its
symbols}\label{lst:b2:cpp-templates:instantiate}

\omoutput{code/bachelor-2/04-cpp-templates/outputs/instantiate.txt}{Thirty instantiated functions, among them
three \code{push_back}s.}\label{out:b2:cpp-templates:instantiate}

\code{nm} lists the symbols of the object file and \code{-C} turns their encoded names back into C++. The three
\code{push_back}s are three separate functions, compiled from one text for three types. Their type \code{W}
(\emph{weak}) tells the linker what to do when several object files contain the same instantiation, which is the
normal case: every file that uses \code{Vec<int>} instantiates it.

\begin{proposition}[One instantiation per argument list]\label{prop:b2:cpp-templates:instantiations}
In each translation unit, a \omterm{def:b2:cpp-templates:template}{template} is instantiated at most once for each distinct list of \omterm{def:b2:cpp-templates:template}{template} arguments,
however many times it is used with that list. When several translation units instantiate the same \omterm{def:b2:cpp-templates:template}{template} with
the same arguments, the linker keeps one copy of each function. A program therefore contains at most one copy of
each instantiated function per distinct argument list, and a \omterm{def:b2:cpp-templates:template}{template} used with $k$ argument lists, of which it
uses $m$ \omterm{def:b2:cpp-classes-raii:class}{member functions} with each, yields at most $km$ functions.
\end{proposition}

\begin{proof}
Within a unit, the compiler records each specialisation it has generated, keyed by the \omterm{def:b2:cpp-templates:template}{template} and its arguments,
and reuses the record for later uses, so a second use generates nothing. Across units, each copy is emitted as a weak
symbol with the same encoded name, built from the \omterm{def:b2:cpp-templates:template}{template} and its arguments; the linker, on finding several weak
definitions of one name, keeps one and discards the others, which the language allows because they are identical
(they come from the same \omterm{def:b2:cpp-templates:template}{template} with the same arguments). The bound follows: at most one function for each pair of
an argument list and a used member.
\end{proof}

\begin{omfigure}
\begin{tikzpicture}[font=\footnotesize, box/.style={draw, rounded corners=2pt, align=center, inner sep=4pt}]
\node[box, fill=omDef!10] (t) at (0,0) {\texttt{template <typename T>}\\ \texttt{class Vec \{\dots\};}};
\node[box, fill=black!4, font=\scriptsize\ttfamily] (u1) at (3.9,1.1) {a.cpp uses\\ Vec<int>, Vec<double>};
\node[box, fill=black!4, font=\scriptsize\ttfamily] (u2) at (3.9,-1.1) {b.cpp uses\\ Vec<int>};
\node[box, fill=omThm!10, font=\scriptsize\ttfamily] (o1) at (8.0,1.1) {a.o (weak symbols):\\ Vec<int>::push\_back\\ Vec<double>::push\_back};
\node[box, fill=omThm!10, font=\scriptsize\ttfamily] (o2) at (8.0,-1.1) {b.o (weak symbol):\\ Vec<int>::push\_back};
\node[box, fill=omMeth!10, font=\scriptsize\ttfamily] (x) at (12.2,0) {program:\\ one Vec<int>::push\_back\\ one Vec<double>::push\_back};
\draw[ptr] (t) -- (u1.west); \draw[ptr] (t) -- (u2.west);
\draw[ptr] (u1) -- (o1); \draw[ptr] (u2) -- (o2);
\draw[ptr] (o1.east) -- (x); \draw[ptr] (o2.east) -- (x);
\foreach \x/\l in {0/template, 3.9/source files, 8.0/compiled: object files, 12.2/linked} \node[font=\scriptsize\bfseries, text=black!60] at (\x,2.3) {\l};
\end{tikzpicture}
\par\smallskip{\small One \omterm{def:b2:cpp-templates:template}{template}, two translation units, two argument lists. Each unit instantiates what it
uses; the linker merges the two copies of \code{Vec<int>::push_back}.}
\end{omfigure}

Two consequences follow. A \omterm{def:b2:cpp-templates:template}{template}'s definition must be visible wherever it is used, since the compiler can only
instantiate what it can read: \omterm{def:b2:cpp-templates:template}{templates} live in headers (\cref{exo:b2:cpp-templates:2} shows what happens
otherwise). And errors in a \omterm{def:b2:cpp-templates:template}{template}'s use appear at instantiation, inside the \omterm{def:b2:cpp-templates:template}{template}'s body, as the next section
shows. Some types deserve a different body from the general one; a \omterm{def:b2:cpp-templates:template}{template} can be given one.

\begin{definition}[Template specialisation]\label{def:b2:cpp-templates:specialisation}
A \emph{full specialisation}\index{template specialisation} \code{template <> ... f<A>(...)} gives the body of a
\omterm{def:b2:cpp-templates:template}{template} for one particular argument list \code{A}, used instead of the general body whenever those arguments are
chosen. \omterm{def:b2:cpp-templates:template}{Function templates} may also be \emph{overloaded}: a second \omterm{def:b2:cpp-templates:template}{template} with a more specific parameter (such as
\code{T *} for \code{T}), chosen when it matches better.
\end{definition}

\omcode{code/bachelor-2/04-cpp-templates/listings/specialise.cpp}{5}{22}{A general \omterm{def:b2:cpp-templates:template}{template}, a specialisation and an
overload}\label{lst:b2:cpp-templates:specialise}

\omoutput{code/bachelor-2/04-cpp-templates/outputs/specialise.txt}{Each value described by the most specific
version.}\label{out:b2:cpp-templates:specialise}

\omterm{def:b2:cpp-templates:template}{Templates} also take values: \code{std::array<double, 3>} is a class with three \code{double}s, its size fixed at
compile time. Values computed at compile time come from \code{constexpr} functions, which the compiler runs while
it compiles when their arguments are constants.

\omcode{code/bachelor-2/04-cpp-templates/listings/constexpr_pow.cpp}{5}{23}{A function run by the compiler, and at
run time}\label{lst:b2:cpp-templates:constexpr}

\omoutput{code/bachelor-2/04-cpp-templates/outputs/constexpr_pow.txt}{Two compile-time values and one run-time
value.}\label{out:b2:cpp-templates:constexpr}

\code{GIB} and the size of \code{table} were computed by the compiler; \code{power(argc + 1, 20)} depends on the
command line, $2^{20}$ here, and was computed when the program ran. A \code{static_assert} is a test that runs at
compile time: if it fails, the program does not compile.

\section{Requirements on a type parameter}\label{sec:b2:cpp-templates:concepts}

A \omterm{def:b2:cpp-templates:template}{template}'s body uses operations of its type parameters (\code{<}, copying, \code{+}), and so requires them. Until
C++20 the requirement was unwritten: the compiler found out at instantiation, inside the body. Below, the same
\omterm{def:b2:cpp-templates:template}{function template} is used with a type that has no \code{<}, first as it stands and then with its requirement written
down.

\omoutput[2-6]{code/bachelor-2/04-cpp-templates/outputs/no_concept.txt}{Without a stated requirement, the error is
reported inside the \omterm{def:b2:cpp-templates:template}{template}.}\label{out:b2:cpp-templates:no-concept}

% ledger: b2:cpp20-concepts
\begin{definition}[Concept, constraint]\label{def:b2:cpp-templates:concept}
A \emph{concept}\index{concept (C++)} is a named, compile-time predicate on \omterm{def:b2:cpp-templates:template}{template} arguments, defined by
\code{template <typename T> concept C = expression;}, where the expression may combine other concepts and
\code{requires} expressions listing operations that must be valid. A \emph{constraint}\index{constraint} attaches
concepts to a \omterm{def:b2:cpp-templates:template}{template}, as \code{template <C T>} or \code{requires C<T>}: the \omterm{def:b2:cpp-templates:template}{template} is then considered only for
arguments that satisfy them. Concepts are a C++20 feature.
\end{definition}

\omcode{code/bachelor-2/04-cpp-templates/listings/with_concept.cpp}{8}{21}{The requirement, stated in
front}\label{lst:b2:cpp-templates:with-concept}

\omoutput[2-5,10-10,15-17]{code/bachelor-2/04-cpp-templates/outputs/with_concept.txt}{With a \omterm{def:b2:cpp-templates:concept}{concept}, the error is
reported at the call, naming the missing operation.}\label{out:b2:cpp-templates:with-concept}

The second message is longer but speaks to the caller: the call does not match, because \code{Point} does not
satisfy \code{Ordered}, because \code{a < b} is invalid. The first speaks of a line inside a function the caller may
never have read. In a library \omterm{def:b2:cpp-templates:template}{template} ten calls deep, the difference is between one message and a page of
them. The standard library predefines many \omterm{def:b2:cpp-templates:concept}{concepts}, such as \code{std::copyable}, \code{std::integral} or
\code{std::random_access_iterator}.

\begin{method}[Stating requirements first]\label{met:b2:cpp-templates:concept-first}
To write a function or \omterm{def:b2:cpp-templates:template}{class template}:
\begin{enumerate}
\item list the operations the body will use on each type parameter, with their result types;
\item express them as a \omterm{def:b2:cpp-templates:concept}{concept}, reusing the library's \omterm{def:b2:cpp-templates:concept}{concepts} where they fit, and constrain the parameter with it;
\item write the body using only those operations, and check that each use is covered (a use the \omterm{def:b2:cpp-templates:concept}{concept} does not
cover compiles for the types tried so far and fails, inside the body, for the next one,
\cref{exo:b2:cpp-templates:6});
\item test the \omterm{def:b2:cpp-templates:template}{template} with two types that satisfy the \omterm{def:b2:cpp-templates:concept}{concept} and one that does not, checking that the last is
refused at the call.
\end{enumerate}
\end{method}

\section{Static against dynamic polymorphism}\label{sec:b2:cpp-templates:static}

A \omterm{def:b2:cpp-templates:template}{function template} and a \omterm{def:b2:cpp-polymorphism:virtual}{virtual function} both let one piece of code work with many types. They choose at different
times.

\begin{definition}[Static polymorphism]\label{def:b2:cpp-templates:static-poly}
\emph{Static polymorphism}\index{static polymorphism} is the use of \omterm{def:b2:cpp-templates:template}{templates} to write code that works with every
type satisfying its requirements, where the operation to run for each type is chosen at compile time, during
instantiation. It is opposed to the \omterm{def:b2:cpp-polymorphism:dispatch}{subtype polymorphism} of \cref{ch:b2:cpp-polymorphism}, chosen at run time
through a \omterm{def:b2:cpp-polymorphism:vtable}{dispatch table}.
\end{definition}

\omterm{def:b2:cpp-templates:static-poly}{Static polymorphism} costs no \omterm{def:b2:cpp-polymorphism:vtable}{table pointer} and no indirect call, and the compiler can inline the chosen operation
into the caller, then optimise the result as a whole. It costs one copy of the code per type, and the set of types
must be known when the program is compiled: a vector of mixed shapes read from a file needs \omterm{def:b2:cpp-polymorphism:dispatch}{dynamic dispatch}. The
two sorts of the standard libraries illustrate the difference. \code{qsort} receives its comparison as a function
pointer and calls it for every comparison; \code{std::sort} is a \omterm{def:b2:cpp-templates:template}{template} whose comparison type is a parameter, so
that each instantiation can inline it.

% ledger: b2:qsort-c89
% ledger: b2:std-sort-complexity
\omcode{code/bachelor-2/04-cpp-templates/listings/sort_compare.cpp}{22}{37}{Sorting the same random integers both
ways; \code{compare_ints} returns the sign of $x - y$}\label{lst:b2:cpp-templates:sort-compare}

\omoutput{code/bachelor-2/04-cpp-templates/outputs/sort_compare.txt}{One run on the author's laptop; the digits
change from run to run.}\label{out:b2:cpp-templates:sort-compare}

\begin{omfigure}
\begin{tikzpicture}[font=\footnotesize, box/.style={draw, rounded corners=2pt, align=center, inner sep=4pt, minimum height=8mm}]
\node[font=\footnotesize\bfseries] at (1.6,1.5) {\texttt{qsort}};
\node[box, fill=black!4] (q) at (0,0) {\texttt{qsort} loop\\ (compiled once)};
\node[box, fill=omThm!10] (c) at (3.4,0) {\texttt{compare\_ints}\\ \scriptsize via a pointer};
\draw[ptr, bend left=25] (q.north east) to node[above, font=\scriptsize] {call, $\approx n\log_2 n$ times} (c.north west);
\draw[ptr, bend left=25] (c.south west) to node[below, font=\scriptsize] {return $-1$, 0 or 1} (q.south east);
\begin{scope}[xshift=7.6cm]
\node[font=\footnotesize\bfseries] at (1.6,1.5) {\texttt{std::sort<int *, Less>}};
\node[box, fill=omMeth!10, minimum width=4.4cm] (s) at (1.6,0) {sort loop with \texttt{x < y} written inside\\ (compiled for \texttt{int} and \texttt{Less})};
\node[font=\scriptsize, text=black!60] at (1.6,-0.9) {no call: one compare instruction};
\end{scope}
\end{tikzpicture}
\par\smallskip{\small Where the comparison happens. \code{qsort} is compiled once for every type, so it can only
call the comparison through a pointer it receives; \code{std::sort} is instantiated for \code{int} and the lambda's
type, and the comparison becomes part of its own loop.}
\end{omfigure}

Both sort the same values in time growing like $n \log n$, and \code{std::sort} is faster by a factor that the run
shows to be about two, from the smallest array to the largest: the function call, the two loads through
\code{void *} and the branchy three-way result of \code{compare_ints}, repeated $n \log_2 n$ times, against a single
comparison inside the loop. The C standard requires nothing of \code{qsort}'s speed: neither its algorithm nor its
complexity is specified. \code{std::sort} must make $O(n \log n)$ comparisons in the worst case.

\begin{inpractice}[Templates or virtual functions]
Use a \omterm{def:b2:cpp-templates:template}{template} when the types are known at compile time and the operation is small and called often: containers,
algorithms, numeric code. Use \omterm{def:b2:cpp-polymorphism:virtual}{virtual functions} when the set of types is open, or chosen at run time from data, or
when one compiled copy must serve every type: plug-ins, the library's fine policies. Mixed designs are common: a
\omterm{def:b2:cpp-templates:template}{template} container of \code{std::unique_ptr} to a \omterm{def:b2:cpp-polymorphism:inheritance}{base class} holds objects of any \omterm{def:b2:cpp-polymorphism:inheritance}{derived class}.
\end{inpractice}

\section{Exercises}\label{sec:b2:cpp-templates:exercises}

\begin{exercise}[$\star$]\label{exo:b2:cpp-templates:1}
For each call a to g, say which \code{show} runs, and with which \code{T} if it is a \omterm{def:b2:cpp-templates:template}{template}. The rule: among the
candidates that match, an exact match beats a conversion, and an ordinary function beats a \omterm{def:b2:cpp-templates:template}{template} when both match
equally well.
\omcode{code/bachelor-2/04-cpp-templates/listings/exo_overload.cpp}{4}{20}{Four candidates, seven calls}
\end{exercise}

\begin{exercise}[$\star$]\label{exo:b2:cpp-templates:2}
A programmer moves the body of a \omterm{def:b2:cpp-templates:template}{function template} into its own source file, as one does for ordinary functions:
\code{twice.hpp} declares \code{template <typename T> T twice(T x);}, \code{twice.cpp} defines it, and
\code{main_twice.cpp} calls \code{twice(21)}. Each file compiles; the link fails. Explain, with
\cref{def:b2:cpp-templates:instantiation}, and give two fixes.
\omoutput[4-6]{code/bachelor-2/04-cpp-templates/outputs/link_twice.txt}{The linker's message.}
\end{exercise}

\begin{exercise}[$\star$]\label{exo:b2:cpp-templates:3}
Write \code{clamp_to(x, lo, hi)}, which returns \code{lo} if $x < lo$, \code{hi} if $hi < x$ and \code{x} otherwise,
for any type with a \code{<}, stated as a \omterm{def:b2:cpp-templates:concept}{concept}; it throws if $hi < lo$. Check that a type without \code{<} is
refused at the call, and say why the function returns a \code{const T &}. Is that safe when \code{lo} is a
temporary?
\end{exercise}

\begin{exercise}[$\star$]\label{exo:b2:cpp-templates:4}
A program uses \code{Vec<int>}, \code{Vec<double>} and \code{Vec<std::string>}, calls \code{push_back} on all three,
\code{at} on the first only, and copies the third. Using \cref{prop:b2:cpp-templates:instantiations}, list the
instantiated \omterm{def:b2:cpp-classes-raii:class}{member functions} of \code{Vec} (count \code{reserve} and \code{clear}, which \code{push_back} and the
\omterm{def:b2:cpp-classes-raii:constructor}{destructor} call). How does the list change if the program is split into two files that both use
\code{Vec<int>}?
\end{exercise}

\begin{exercise}[$\star\star$]\label{exo:b2:cpp-templates:5}
Two of these three calls do not compile; the third compiles and is wrong. Explain the two messages with
\cref{def:b2:cpp-templates:deduction}, fix each call, and explain what the third compares.
\omcode{code/bachelor-2/04-cpp-templates/listings/exo_deduce.cpp}{5}{14}{Three calls}
\omoutput[3-3,10-10,13-13,20-20]{code/bachelor-2/04-cpp-templates/outputs/exo_deduce.txt}{The key lines of the two
errors.}
\end{exercise}

\begin{exercise}[$\star\star$]\label{exo:b2:cpp-templates:6}
The \omterm{def:b2:cpp-templates:concept}{concept} below is satisfied by \code{double} and by \code{std::string}, and yet the call with strings does not
compile. Explain the message, and say which step of \cref{met:b2:cpp-templates:concept-first} was skipped. Write the
\omterm{def:b2:cpp-templates:concept}{concept} the body really needs.
\omcode{code/bachelor-2/04-cpp-templates/listings/exo_weak_concept.cpp}{7}{28}{A middle value}
\omoutput[2-6]{code/bachelor-2/04-cpp-templates/outputs/exo_weak_concept.txt}{The error, inside the body.}
\end{exercise}

\begin{exercise}[$\star\star$]\label{exo:b2:cpp-templates:7}
Write a \omterm{def:b2:cpp-templates:template}{function template} \code{lower_bound_of(first, last, key)} that returns the first position $p$ of a sorted
range with \code{!(*p < key)}, or \code{last}, for any random-access iterator (a \omterm{def:b2:cpp-templates:concept}{concept} of the library) and any key
comparable with the elements. State its loop invariant. Test it against \code{std::lower_bound} on random sorted
vectors and every key in a range.
\end{exercise}

\begin{exercise}[$\star\star$]\label{exo:b2:cpp-templates:8}
A \omterm{def:b2:cpp-templates:template}{class template} \code{Pow<N>}, with a \omterm{def:b2:cpp-templates:template}{template parameter} \code{unsigned N}, has one member, the constant
\code{static constexpr long value = 3 * Pow<N - 1>::value;}, and a \omterm{def:b2:cpp-templates:specialisation}{full specialisation} \code{Pow<0>} whose
\code{value} is 1: \code{Pow<N>::value} is $3^N$, computed at compile time. How many classes does
\code{Pow<20>::value} instantiate? And a program that uses \code{Pow<20>} and \code{Pow<25>}? Write the same with a
\code{constexpr} function and compare what the compiler must generate.
\end{exercise}

\begin{exercise}[$\star\star$]\label{exo:b2:cpp-templates:9}
Write \code{Matrix<T, R, C>}, stored in a \code{std::array<T, R * C>}, with element access, \code{transpose} and a
product \code{operator*} that accepts only matrices of matching inner dimension, so that multiplying a $2 \times 3$
by a $2 \times 3$ matrix does not compile. Check that refusal with a \code{static_assert} and a \omterm{def:b2:cpp-templates:concept}{concept}.
\end{exercise}

\begin{exercise}[$\star\star\star$]\label{exo:b2:cpp-templates:10}
Write a binary heap \code{Heap<T, Less = std::less<T>>} over a \code{std::vector<T>}, with \code{push}, \code{top} and
\code{pop}, the comparator being a \omterm{def:b2:cpp-templates:template}{template parameter} stored in the object. Test it against a sorted copy for 500
random values, and as a min-heap of strings with \code{std::greater}. Why can the comparator be inlined here and not
in \code{qsort}?
\end{exercise}

\begin{exercise}[$\star\star\star$]\label{exo:b2:cpp-templates:11}
A library \omterm{def:b2:cpp-templates:template}{template} has $m$ \omterm{def:b2:cpp-classes-raii:class}{member functions} of about $b$ bytes of machine code each, and a program uses it with $k$
types, every member with every type. (a) Bound the code size the \omterm{def:b2:cpp-templates:template}{template} contributes, with
\cref{prop:b2:cpp-templates:instantiations}. (b) The program is built from $u$ files that each use all $k$ types:
what do the object files contain before linking, and the program after? (c) From
\cref{out:b2:cpp-templates:instantiate}, how many functions of \code{Vec} does \code{vec_demo.o} contain per type,
roughly? Compare with a \code{void *} version, compiled once.
\end{exercise}

\begin{exercise}[$\star\star\star$]\label{exo:b2:cpp-templates:12}
Write \code{pascal<N>()}, a \code{constexpr} function returning the table of binomial coefficients $\binom{n}{k}$ for
$0 \leq k \leq n \leq N$ as nested \code{std::array}s, computed by Pascal's rule. Keep it in a \code{constexpr}
variable for $N = 67$, and check with \code{static_assert} that $\binom{5}{2} = 10$ and that $\binom{67}{33}$ fits in 64
bits. Why is 67 the largest $N$ for which every entry fits in an unsigned 64-bit integer? Test Pascal's rule and the
row sums $2^n$ at run time.
\end{exercise}

\section{Problem: One sort for every type}\label{sec:b2:cpp-templates:problem}

\begin{problem}[{Weekend problem --- One sort for every type}]\label{pb:b2:cpp-templates:1}
The town keeps tables of readings (\code{double}), of people sorted by name (\code{std::string}) and of loans sorted
by date (a \code{struct}). This problem sorts all three with one generic function and measures what it gains over
\code{qsort}.

\medskip\noindent\textbf{Part I --- \code{qsort} and \code{void *}.}\par\nopagebreak
\begin{enumerate}[beginpenalty=10000]
\item Write the comparison function \code{qsort} needs for \code{double}s. Why is \code{return x - y;} wrong for
\code{double}, and for \code{int} near the ends of its range?
\item Write the comparison for an array of \code{struct loan {int id; int due;}} by due date, then by
\code{id}.
\item Why can \code{qsort} not sort an array of \code{std::string} correctly? (It moves elements byte by byte.)
\item What does the compiler check about the arguments of a \code{qsort} call? Give a call that compiles and sorts
garbage.
\item \code{compare_ints} in \cref{lst:b2:cpp-templates:sort-compare} is called how many times, roughly, for
$n = 10^6$? Use $n \log_2 n$.
\item What does the C standard promise about \code{qsort}'s running time and stability?
\end{enumerate}

\medskip\noindent\textbf{Part II --- A template sort.}\par\nopagebreak
\begin{enumerate}[resume, beginpenalty=10000]
\item Write \code{template <typename T, typename Less> void insertion_sort(T *a, std::size_t n, Less less)}.
\item Call it on \code{double}s with a lambda, on \code{std::string}s with \code{std::less<>}, and on loans with a
lambda comparing due dates then ids.
\item Which functions does the compiler instantiate for these three calls?
\item Why are the strings now moved correctly, where \code{qsort} copied bytes?
\item Rewrite the swap of two elements so that it moves, and say what it costs for strings longer than the small
buffer of \cref{ch:b2:cpp-copy-move}.
\item Insertion sort makes up to $n(n-1)/2$ comparisons. Why would a \omterm{def:b2:cpp-templates:template}{template} version still lose to \code{qsort} for
$n = 10^6$?
\end{enumerate}

\medskip\noindent\textbf{Part III --- Concepts.}\par\nopagebreak
\begin{enumerate}[resume, beginpenalty=10000]
\item State, as a \omterm{def:b2:cpp-templates:concept}{concept}, what \code{insertion_sort} requires of \code{Less} given \code{T}.
\item Which \omterm{def:b2:cpp-templates:concept}{concept} of the library expresses ``\code{T} can be moved and move-assigned''? Add it.
\item A caller passes a comparator that takes its arguments by non-const reference. Does it satisfy your \omterm{def:b2:cpp-templates:concept}{concept}?
Should it?
\item A comparator \code{[](double x, double y) {return x <= y;}} satisfies the \omterm{def:b2:cpp-templates:concept}{concept}. Why is it still wrong for a
sort? Can a \omterm{def:b2:cpp-templates:concept}{concept} detect it?
\item Rewrite the sort to take a pair of random-access iterators instead of a pointer and a size. Which library
\omterm{def:b2:cpp-templates:concept}{concept} do you use?
\item Write a test that checks, for each of the three types, that the result is sorted and is a permutation of the
input.
\end{enumerate}

\medskip\noindent\textbf{Part IV --- The measured speed-up.}\par\nopagebreak
\begin{enumerate}[resume, beginpenalty=10000]
\item From \cref{out:b2:cpp-templates:sort-compare}, compute the time per element and per comparison of each sort
for $n = 10^6$.
\item Is the ratio roughly constant over $n$? What does that say about where the difference comes from?
\item Explain, with the figure on the two sorts, what the processor does per comparison in each case.
\item The town sorts a million readings once a night. Is the difference worth a rewrite? And a trading system that
sorts ten thousand orders every millisecond?
\item The \code{sort_compare} program was compiled with \code{-O2}. Predict the ratio at \code{-O0}, where nothing
is inlined.
\item State the result of the problem: the speed-up of \code{std::sort} over \code{qsort} measured on the laptop, and
its cause.
\end{enumerate}
\end{problem}
